\section{Introduction}\label{sec:intro}

For $n\ge 2$, $a>0$, $\lambda\ge 0$ and $t>0$ let
\[
R^\lambda_{a,t}f(x)=\frac{1}{(2\pi)^n}\int_{\R^n}\Bigl(1-\frac{|\xi|^a}{t^a}\Bigr)_+^\lambda
\hat f(\xi)\,e^{i\ip{x}{\xi}}\,d\xi ,
\]
where $\hat f(\xi)=\int_{\R^n}f(y)e^{-i\ip{y}{\xi}}\,dy$.
For $a=2$ these are the Bochner--Riesz means of the Fourier inversion integral, and for
$a=1$ the Riesz means. For $1\le p<2n/(n+1)$ the number
\begin{equation}\label{eq:critical}
\lambda(p)=n\Bigl(\frac1p-\frac12\Bigr)-\frac12
\end{equation}
is the critical index: $R^\lambda_{a,t}$ is unbounded on $L^p(\R^n)$ for $\lambda\le\lambda(p)$,
because its kernel decays like $|x|^{-(n+1)/2-\lambda}$ and fails to lie in $L^p$ when
$\lambda\le\lambda(p)$ \cite{Ste93}. The Bochner--Riesz conjecture asserts boundedness on $L^p$ for
every $\lambda>\lambda(p)$ when $1\le p<2n/(n+1)$. At the critical index itself the
expected statement is the \emph{weak-type endpoint Bochner--Riesz conjecture}: for
$1<p<2n/(n+1)$, the operators $R^{\lambda(p)}_{a,t}$ map $L^p(\R^n)$ to the Lorentz space
$L^{p,\infty}(\R^n)$, uniformly in $t$. This is the title problem of Tao's paper \cite{Tao98}.

The case $p=1$, with index $\lambda(1)=(n-1)/2$, is due to Christ \cite{Chr88}, and was
extended by Vargas to $L^1(w)\to L^{1,\infty}(w)$ for every $A_1$ weight $w$ \cite{Var96}.
For $p>1$, Christ proved the weak-type endpoint in the range $1<p<2(n+1)/(n+3)$
\cite{Chr87}, Tao proved it at $p=2(n+1)/(n+3)$ \cite{Tao98}, and Seeger settled the plane \cite{See96}.
Tao then showed that for $1<p<2n/(n+1)$ the weak-type endpoint follows from the strong-type
Bochner--Riesz bounds for all $\lambda>\lambda(p)$, by an $\varepsilon$-removal argument
\cite{Tao98}. In higher dimensions only partial ranges were known, since the Bochner--Riesz
conjecture itself was open for $n\ge 3$.

On 25 October 2026 OpenAI released a proof of Sogge's local smoothing conjecture for the
Euclidean wave equation in every dimension \cite{OAI26}, and deduced from it the
Bochner--Riesz conjecture for every $n\ge 2$ and every $1\le p\le\infty$. Combined with the
conditional theorems of Tao \cite{Tao98} and of Beltran, Roos and Seeger \cite{BRS25,BRS23},
this gives the endpoint bounds below in every dimension. The implications themselves are not
new: Tao showed that the weak-type endpoint follows from the Bochner--Riesz conjecture
\cite{Tao98}, and Beltran, Roos and Seeger state the weighted range of Theorem C, conditionally on
that conjecture, in the introduction to \cite{BRS25}. What this paper adds is a complete account,
in one place, of the argument from local smoothing to the endpoint bounds: the two inputs that
have not been refereed are reproduced in full, with credit, and the conditional theorems are
checked against them. Consequences of another OpenAI release, for weighted counting and exact
optimisation, are derived in the same manner in the paper archived as \cite{Bha26b}.

\subsection{Results}
Write $\Sp(p,q)$ for the class of operators with a $(p,q)$ sparse bound (Definition
\ref{def:sparse}); $A_1$ and $RH_\sigma$ are the Muckenhoupt and reverse H\"older classes
(Section \ref{sec:weights}).

\begin{theoremA}[Weak-type endpoint]
Let $n\ge2$, $a>0$ and $1<p<2n/(n+1)$. Then
\[
\sup_{t>0}\bigl\|R^{\lambda(p)}_{a,t}f\bigr\|_{L^{p,\infty}(\R^n)}\le C_{n,a,p}\|f\|_{L^p(\R^n)},
\]
and $\|R^{\lambda(p)}_{a,t}f-f\|_{L^{p,\infty}}\to0$ as $t\to\infty$ for every $f\in L^p(\R^n)$.
\end{theoremA}

\begin{theoremB}[Sparse bounds]
Let $n\ge2$, $a>0$, $1\le p<2n/(n+1)$ and
\[
q_{\rm opt}(p)=\frac{(n-1)p}{n+1-2p}<q<p'.
\]
Then $R^{\lambda(p)}_{a,t}\in\Sp(p,q)$, with norm independent of $t$.
\end{theoremB}

The upper bound $q<p'$ only excludes trivial cases: for $1/p+1/q\le1$ every nonzero pair of
compactly supported functions has infinite sparse form (Lemma \ref{lem:trivialsparse}). For
$\lambda=\lambda(p)$ the pairs $(1/p,1/q)$ of Theorem B fill the critical vertical side $P_1P_2$ of
the trapezoid of exponents described in \cite[(1.5)]{BRS25}, except for its upper corner $P_2$,
where $q=q_{\rm opt}(p)$. Sparse bounds fail outside the closed trapezoid \cite{BBL17}, so on the
critical side Theorem B is sharp except at $P_2$. The corner is known at the Stein--Tomas exponent
$p=2(n+1)/(n+3)$ \cite[Theorem 1.5]{BRS25} and is open otherwise; Section \ref{sec:open} explains
why the method stops short of it.

\begin{theoremC}[Weights]
Let $n\ge 2$, $a>0$, and let $w\in A_1$ belong to $RH_{\sigma}$ for some $\sigma>1$. Put
\[
p_1(w)=1+\frac{n-1}{n+1}\Bigl(1-\frac1{\sigma}\Bigr).
\]
Then for $1<p<p_1(w)$ the operators $R^{\lambda(p)}_{a,t}$ map $L^p(w)$ to $L^{p,\infty}(w)$
uniformly in $t$, and $R^{\lambda(p)}_{a,t}f\to f$ in $L^{p,\infty}(w)$ for every $f\in L^p(w)$.
\end{theoremC}

Every $A_1$ weight lies in some $RH_\sigma$ with $\sigma>1$, so Theorem C gives a range
$p>1$ for every $A_1$ weight, as in \cite[Theorem 1.1]{BRS25}. Unconditionally, \cite[\S9]{BRS25}
caps the reverse H\"older exponent at $(n+3)/2$, which limits the range to $p<2(n+1)/(n+3)$; the
uncapped range is stated there as a consequence of the Bochner--Riesz conjecture
\cite[\S1.1]{BRS25}, and Theorem C is that statement. As $\sigma\to\infty$ the range fills the
whole interval $1<p<2n/(n+1)$, and the constant weight recovers Theorem A. The range is also
open at its right end: by the self-improvement of reverse H\"older inequalities, every
$w\in A_1\cap RH_\sigma$ lies in $RH_\Sigma$ for an explicit $\Sigma>\sigma$, so the conclusion of
Theorem C holds at $p=p_1(w)$ as well (Theorem \ref{thm:boundary}). For $w\in A_1$ alone it holds
for
\[
1<p<1+\frac{n-1}{(n+1)\bigl(2^{n+1}\kappa_n[w]_{A_1}+1\bigr)},
\]
with $\kappa_n$ the ratio of the volume of a ball to that of its inscribed cube (Corollary
\ref{cor:A1only}).

For $n=2$ the Bochner--Riesz conjecture is due to Carleson and Sj\"olin \cite{CS72}, and Theorems A--C
were proved in \cite{See96,BRS25}. For $n\ge3$ they follow from \cite{OAI26} and the conditional
theorems above.

\subsection{What is borrowed and what is reproduced}
Theorems A--C are obtained from three ingredients.
\begin{enumerate}[label=(\roman*)]
\item The Bochner--Riesz conjecture, proved by OpenAI \cite{OAI26} through local smoothing.
\item The off-diagonal $\varepsilon$-removal of Tao \cite{Tao98}, in the vector-valued form of
Beltran, Roos and Seeger \cite{BRS23}: non-endpoint bounds imply the multi-scale
inequalities $\VBR(p,r)$ of Definition \ref{def:vbr}.
\item The sparse domination theorem of Beltran, Roos and Seeger \cite[Theorem 2.3]{BRS25}, which
turns $\VBR(p,r)$ into sparse bounds, and the weighted theorem of Frey and Nieraeth
\cite[Theorem 1.4]{FN19}, which turns sparse bounds into weighted weak-type bounds.
\end{enumerate}
Items (iii) are refereed and are cited. Items (i) and (ii) are a preprint and an Oberwolfach
preprint, and their arguments are reproduced here in full, with credit: Part I gives OpenAI's proof of the
critical annular estimate, of local smoothing and of the Bochner--Riesz conjecture, with the
Dirichlet problem it uses in Appendix \ref{app:dirichlet}; Part II gives the
$\varepsilon$-removal. The arguments are those of \cite{OAI26} and \cite{BRS23}; we follow
them closely, specialise them to the round sphere, which is all that is needed, and supply
some intermediate computations. Part III proves Theorems A--C.

Reproducing these arguments is not an independent proof of local smoothing: the analysis of
Part I is that of \cite{OAI26}, and its correctness is what Theorems A--C rest on in dimensions
$n\ge3$. The implications from the Bochner--Riesz conjecture to Theorems A--C are due to Tao and
to Beltran, Roos and Seeger, as described above, and we present Theorems A--C as their
consequences.

\subsection{Organisation}
Sections \ref{sec:current}--\ref{sec:euclid} form Part I. Section \ref{sec:current} builds a
positive conserved current from a wave; Section \ref{sec:hessian} develops the quaternionic
Hessian and its cofactor; Section \ref{sec:energy} defines a determinant energy on spatial
fibres; Section \ref{sec:fiber} estimates it; Section \ref{sec:cylinder} compares two circles in
the auxiliary time plane; Section \ref{sec:euclid} removes the period and proves local
smoothing and the Bochner--Riesz conjecture. Section \ref{sec:vbr} is Part II. Section
\ref{sec:endpoint} proves Theorems A--C, Section \ref{sec:boundary} shows that the range of the
weights is open at its right end, and Section \ref{sec:open} lists what remains open.
Exact-arithmetic checks of the algebraic identities and of the exponent computations are
available with the sources \cite{Bha26arch}; no statement depends on them.

\subsection{Notation}
In Part I the spatial dimension is $n\ge2$ and $d=n-1$; in Parts II and III the dimension is
again $n$, and $d$ is not used. We write $A\lesssim B$ when $A\le CB$ with a constant
depending only on fixed parameters, and $A\asymp B$ when both $A\lesssim B$ and $B\lesssim A$.
The Fourier multiplier with symbol $m$ is $m(D)$.
